\documentclass[11pt,a4paper]{article}

\usepackage[T1]{fontenc}
\usepackage{lmodern}
\usepackage[margin=25mm]{geometry}
\usepackage{microtype}
\usepackage{amsmath,amssymb}
\usepackage{graphicx}
\usepackage{booktabs,tabularx}
\usepackage{xcolor}
\usepackage{float}
\usepackage{caption}
\usepackage{fancyhdr}
\usepackage[colorlinks=true,linkcolor=ink,urlcolor=accent,citecolor=ink]{hyperref}

\graphicspath{{figures/}{../images/views/}{../images/}{../images/sketch/stills/}}

\definecolor{ink}{HTML}{1B2A41}
\definecolor{accent}{HTML}{C0394B}
\definecolor{boxbg}{HTML}{F3F6F9}
\definecolor{boxline}{HTML}{9AA5B1}

% consistent look: sans-serif headings, one callout style, one figure style
\makeatletter
\renewcommand\section{\@startsection{section}{1}{0pt}{-3ex plus -1ex minus -.2ex}{1ex plus .2ex}{\sffamily\Large\bfseries\color{ink}}}
\renewcommand\subsection{\@startsection{subsection}{2}{0pt}{-2.4ex plus -1ex minus -.2ex}{0.6ex plus .2ex}{\sffamily\large\bfseries\color{ink}}}
\makeatother
\newsavebox{\calloutbox}
\newenvironment{callout}[1]{%
  \par\medskip\noindent\begin{lrbox}{\calloutbox}\begin{minipage}{\dimexpr\linewidth-2\fboxsep-2\fboxrule\relax}%
  \small\textsf{\textbf{\textcolor{accent}{#1}}}\par\smallskip}%
  {\end{minipage}\end{lrbox}\fcolorbox{boxline}{boxbg}{\usebox{\calloutbox}}\par\medskip}
\newcommand{\fig}[4]{\begin{figure}[!htb]\centering\includegraphics[width=#2\linewidth]{#1}\caption{#3}\label{#4}\end{figure}}

\pagestyle{fancy}\fancyhf{}\renewcommand{\headrulewidth}{0pt}\fancyfoot[C]{\small\thepage}
\setlength{\parindent}{0pt}\setlength{\parskip}{0.6em}
\captionsetup{font=small,labelfont={bf,sf},format=hang,skip=6pt}
\setcounter{secnumdepth}{2}
\renewcommand{\arraystretch}{1.15}

\begin{document}

{\sffamily\color{ink}\Large\bfseries A 3D-printed planimeter\par}
{\sffamily\color{ink}\large How it measures area, the maths behind it, and how the maths became plastic\par}
\vspace{0.4em}\hrule\vspace{0.8em}

\fig{real.png}{0.7}{The finished design: a pole weight on the left, a tracer arm at the front, and a small wheel and counter at the right.}{fig:real}

\section{What a planimeter does}

A planimeter measures the area of a shape. You rest a pointer on the edge of the shape, follow the edge all the way
round until you are back where you started, and read a number off a small wheel. That number is the area inside the
outline. It works for a lake on a map, a leaf, or a cross section of a ship's hull, and the outline can be as wobbly as
you like.

This article explains \emph{why} that works. It assumes you know one thing: that an integral adds up lots of small
pieces. Everything else is built from scratch, and then matched to the parts of a design you can print at home.

\begin{callout}{The whole idea in three sentences}
A wheel only rolls when you push it sideways to its axle. A rod that sweeps over a patch of ground pushes its wheel
sideways by an amount that depends on the area it sweeps. Go round a closed outline, and the wheel has counted that area
for you.
\end{callout}

\section{The maths, from scratch}

\subsection{Adding up small pieces: the integral}

To find the area under a curve, cut it into thin vertical strips. Each strip is almost a rectangle, so its area is
height times width. Add the strips up. The thinner the strips, the closer the sum is to the true area, and the exact
answer is written $\int f(x)\,dx$. The stretched ``S'' stands for ``sum''.

\fig{fig_strips.pdf}{0.95}{An integral is a sum of thin strips.}{fig:fig_strips}

\subsection{Adding up along a path: the path integral}

Now suppose you are not adding along a straight line but along a curved path, such as a walk through a park. Cut the
path into tiny steps. If you add up the \emph{lengths} of the steps, you get the length of the whole path. That is the
simplest path integral.

A wheel adds up something different. A wheel on a trolley rolls when you push it forward but just skids when you push
it along its axle. For each tiny step, only the part at right angles to the axle makes the wheel turn. The wheel's
total reading is the sum of those counted parts along the path. When the path is a closed loop that ends where it
began, the sum is written with a loop on the integral sign, $\oint$.

\fig{fig_pathsum.pdf}{0.95}{Left: a path integral adds up small steps along a path. Right: a wheel counts only the part of each step that is at right angles to its axle.}{fig:fig_pathsum}

\subsection{Meet the instrument}

The planimeter has two arms (Figure~\ref{fig:linkage}). The pole arm is fixed at one end to a heavy pole $O$, so its
other end, the elbow $E$, can only swing on a circle. The tracer arm is a rod of length $L$ joined to the elbow, and the
pointer $T$ sits at its far end. A small wheel is fixed to the tracer arm with its axle parallel to the arm, at a
distance $a$ from the elbow.

\fig{fig_linkage.pdf}{0.8}{The two arms. The wheel's axle is parallel to the tracer arm, so the wheel rolls sideways to it.}{fig:linkage}

\begin{table}[H]
\centering\small
\begin{tabular}{@{}cl@{\qquad}cl@{}}
\toprule
Symbol & Meaning & Symbol & Meaning \\
\midrule
$L$ & length of the tracer arm & $a$ & wheel position along the arm \\
$R$ & length of the pole arm & $s$ & total distance the wheel rolled \\
$A$ & area of the outline & $\oint$ & sum around a closed loop \\
\bottomrule
\end{tabular}
\caption{Symbols used in this article.}
\end{table}

\subsection{The area swept by a rod}

Here is the key step. Picture the tracer arm as a rod that moves a little. Some points on it move further sideways than
others, because turning moves the far end more than the near end. The sideways push grows steadily from one end to the
other, so on a graph it is a straight line (Figure~\ref{fig:rod}).

The area swept is the sum, along the rod, of the sideways push at each point. That is the area under the straight line,
and the area under a straight line equals its length times its value at the middle. So:
\begin{equation}
\text{area swept by the rod} \;=\; L \times (\text{sideways push of the rod's middle}).
\end{equation}

\fig{fig_rod.pdf}{0.95}{Left: a rod slides and turns, sweeping the shaded area. Right: the sideways push varies in a straight line along the rod, so the area is the length times the push at the middle.}{fig:rod}

A wheel at the middle of the rod would count exactly that push. A wheel at some other position $a$ counts the push at
$a$ instead, which is slightly different. The next step shows that this difference disappears over a full loop.

\subsection{Going once round the outline}

As the pointer follows the outline (Figure~\ref{fig:loop}), the tracer arm turns, but it ends the loop pointing the
same way it started. Its total turning is zero. The push at any position along the rod equals the push at the middle
plus an amount proportional to how fast the rod is turning. Added up over the whole loop, the turning part adds up to
zero. So a wheel at \emph{any} position reads the same total, $s$, and that total is the one a wheel at the middle would read.

The elbow, meanwhile, only rocks to and fro along one arc and comes back, so the path it traces encloses no area. All
the ground swept by the rod is therefore the inside of the pointer's loop:
\begin{equation}
\boxed{\;A \;=\; L \times s\;}
\end{equation}

\fig{fig_loop.pdf}{0.8}{The pointer follows the outline. The elbow rocks along an arc and comes back, so only the inside of the outline is counted.}{fig:loop}

\begin{callout}{What this means for the design}
The area is the arm length times the wheel's reading, and nothing else. The wheel can sit anywhere along the arm. The
length of the pole arm does not matter. Only $L$ and the size of the wheel need to be right, and both can be calibrated
(Section~\ref{sec:use}).
\end{callout}

\subsection{If the pole is inside the outline}

If the pole is inside the outline, the elbow goes once round its whole circle and the rod turns through one full turn
($2\pi$). Repeating the argument adds a fixed amount:
\begin{equation}
A \;=\; L\,s \;+\; \pi\,(R^2 + L^2 - 2aL).
\end{equation}
The extra term is the area the instrument shows when the wheel reads zero. For this design ($R=160$\,mm, $L=100$\,mm,
$a=-22$\,mm) it is $\pi\times 200^2\,\text{mm}^2\approx1257$\,cm$^2$. It is easier to keep the pole outside the outline,
so the design is made for that.

\subsection{Checking the algebra with a computer}

A sign slip is easy to make, so the formula was tested numerically. A script moves the pointer around known shapes,
finds where the elbow must be at every step, rolls the wheel, and compares $L\times s$ with the exact area.

\begin{table}[H]
\centering\small
\begin{tabular}{llrrr}
\toprule
Pole & Shape & Exact (cm$^2$) & Simulated (cm$^2$) & Error \\
\midrule
outside & circle, radius 40\,mm & 50.265 & 50.265 & $<0.0001\%$ \\
outside & square, 80\,mm & 64.000 & 64.000 & $<0.0001\%$ \\
outside & ten-point star & 26.450 & 26.450 & $<0.0001\%$ \\
inside & circle, radius 150\,mm & 706.849 & 706.851 & 0.0002\% \\
\bottomrule
\end{tabular}
\caption{The formula against exact areas. The tiny differences come from the finite step size.}
\end{table}

\section{From maths to plastic}

Each quantity in the formula became a physical feature (Table~\ref{tab:map}).

\begin{table}[H]
\centering\small
\begin{tabularx}{\linewidth}{@{}llX@{}}
\toprule
In the maths & On the instrument & Value \\
\midrule
Arm length $L$ & Tracer arm slides in a channel, set on a scale & 100\,mm (adjustable 90 to 110) \\
Wheel size & Rolling circumference of the knurled wheel & 100\,mm (diameter 31.83) \\
Reading $s$ & Drum around the wheel, plus a vernier & 0.2\,cm$^2$ steps \\
Number of turns & Ten-slot Geneva counter & 0 to 9 turns \\
Pole arm $R$ & Bar with a bearing at each end & 160\,mm \\
Wheel position $a$ & Wheel mounted beside the arm & $-22$\,mm \\
Wheel axle parallel to arm & Steel axle through two printed plates & within $0.2^\circ$ (target) \\
\bottomrule
\end{tabularx}
\caption{From the formula to the hardware.}\label{tab:map}
\end{table}

\subsection{Round numbers}

The arm is 100\,mm and the wheel's circumference is 100\,mm, so one turn of the wheel is $100\times100=10\,000$\,mm$^2$,
which is 100\,cm$^2$. One millimetre of wheel travel is exactly one square centimetre. The wheel's diameter is
$100/\pi=31.83$\,mm.

\subsection{Reading the wheel}

The drum around the wheel has 50 marks, so each mark is $100/50=2$\,cm$^2$. A fixed vernier next to the drum gives the
tenths of a mark. It has ten divisions that together span nine drum marks, so each is a little shorter than a drum mark.
Look for the vernier line that lines up with a drum line. Its number times 0.2 is added to the drum reading
(Figure~\ref{fig:reading}).

\fig{fig_reading.pdf}{0.95}{Reading the drum and vernier. The vernier zero is just past 36 and line 3 lines up, so the reading is $36 + 3\times0.2 = 36.6$\,cm$^2$.}{fig:reading}

\subsection{Counting full turns}

One turn of the wheel is 100\,cm$^2$, so a larger area needs several turns. A small pin on the wheel flicks a toothed
dial forward one slot each turn, the same trick a film projector uses. This is called a Geneva drive. The dial has ten
slots, so one turn of the wheel moves it $360^\circ/10=36^\circ$, and the dial shows the digits 0 to 9. A raised disc
on the wheel locks the dial still between flicks. The pin radius $a_p$ fixes the spacing between the wheel and the dial
axles, $C=a_p/\sin(\pi/10)=19.4$\,mm for $a_p=6$\,mm.

\fig{fig_geneva.pdf}{0.8}{Dial angle against wheel angle over one turn. The dial moves only while the pin is engaged (shaded) and is locked the rest of the time.}{fig:fig_geneva}

The timing is set so that the dial completes its step just as the drum passes from 99 to 0. Between digits, read the
lower one.

\subsection{The printed parts}

The design has nine printed parts, plus a calibration rule and a test coupon. Figure~\ref{fig:exploded} shows how
everything fits together.

\fig{exploded.png}{1.0}{Exploded view of every part and fastener.}{fig:exploded}

\begin{table}[H]
\centering\small
\begin{tabularx}{\linewidth}{@{}llX@{}}
\toprule
Part & Printed & Job \\
\midrule
Wheel & drum end down & Knurled rim, graduated drum, bearing pockets, Geneva pin \\
Block A & on end & Channel that holds the tracer arm; elbow post and clamp \\
Plates B and C & flat & Wheel mounts; B carries the vernier and a ballast box, C the dial axle \\
Counter dial & flat & Ten-slot Geneva wheel with digits \\
Tracer arm & flat & $12\times8$\,mm bar with pivot slot, scale and knob \\
Pole arm & flat & Bar with a bearing at each end \\
Pole weight and lid & upright & Keeps the pole still; holds ballast \\
Check rule, coupon & flat & Calibration circle of 100\,cm$^2$; test print for fits \\
\bottomrule
\end{tabularx}
\caption{The printed parts. Orientations avoid supports.}
\end{table}

Three design choices come straight from the maths. The elbow bolt passes through the centreline of the tracer arm,
because a sideways offset would add a term the wheel does not count. The wheel runs on a steel axle held by two
bearings, so the printed plastic does not decide the direction of the axle. The carriage is split into three printed
pieces so that every piece can lie flat or stand on end without supports.

\section{Does it hold up?}

\subsection{Checks on the model}

\begin{itemize}
\item All 42 components sit in the working pose with no overlaps.
\item The pole arm swings through every elbow angle from $8^\circ$ to $179^\circ$ without touching the carriage, wheel or dial.
\item The Geneva dial turns by exactly $36^\circ$ per wheel turn, with no collisions at any step.
\item With a few grams of steel shot in the ballast box, about 70\,g presses the wheel on the paper, enough to stop it
slipping.
\end{itemize}

\subsection{Stiffness and strength}

Under a harsh 2\,N push on the pointer (about 200\,g), the tracer arm bends 0.2\,mm down and 0.09\,mm sideways, and its
stress stays 22 times below the plastic's limit. The pole arm bends 0.1\,mm under 1\,N at mid-span. The one tight spot
is the press fit of the bearings: at 0.05\,mm of interference the margin is only $1.4\times$, so the coupon exists to
find a gentle fit.

\subsection{How accurate can it be?}

Two small flaws matter most: a wheel axle that is not parallel to the arm, and play in the elbow bearing. Simulating
both gives Figure~\ref{fig:err}. An error in the wheel diameter or the arm length only scales every reading, and the
check rule removes it. A realistic target is about 0.5\,\% for shapes of 50 to 300\,cm$^2$, a prediction until a
printed copy is measured.

\fig{fig_errors.pdf}{0.95}{Worst-case area error from a misaligned wheel axle (left) and elbow play (right).}{fig:err}

One limit is built in. Professional planimeters average two readings taken with the pole on either side to cancel
axle misalignment. In this design the pole arm would hit the wheel on one side, so only one side is possible.

\section{Build, calibrate, measure}\label{sec:use}

\textbf{Build.} Print the coupon first and adjust the three fit values in \texttt{cad/params.py} if needed. Then print
the parts, press the bearings in, and assemble using the exploded view and the steps in the repository README.

\textbf{Calibrate.} Pin the check rule to a board and rest the pointer in the cone seat. The circle it traces has area
exactly 100\,cm$^2$. Trace it once with the pole outside. If the instrument reads $R_\text{meas}$, set the arm length to
\begin{equation}
L_\text{new} = L_\text{old}\times\frac{100}{R_\text{meas}},
\end{equation}
and repeat until it reads 100.

\textbf{Measure.} Put the pole weight outside the outline. Rest the pointer on the outline, note the dial and drum,
trace clockwise back to the start, and note them again. The difference is the area in square centimetres.

\section{What is still unproven}

The model, the checks and the simulation agree with each other, but only a printed copy can settle some things.
Printed holes come out slightly small, so the bearing and arm fits depend on the coupon. The wheel load is worked out
from an assumed infill and should be weighed. A bare PLA rim might slip on smooth paper, and an O-ring fixes that.
Finally, check that the wheel rolls straight along the arm before trusting small shapes.

\medskip
\noindent\small Source files, STL and STEP models, and the simulation are in the project repository.

\end{document}
